\documentclass[11pt]{article}
\usepackage[a4paper,margin=28mm]{geometry}
\usepackage[T1]{fontenc}
\usepackage{lmodern}
\usepackage{amsmath,amssymb,amsthm,mathtools}
\usepackage{booktabs,longtable,array}
\usepackage{microtype}
\usepackage[hidelinks]{hyperref}
\usepackage{url}
\hypersetup{
 pdftitle={Kusner's Conjecture in l1^5},
 pdfauthor={Xiao Li}
}
\allowdisplaybreaks
\newtheorem{theorem}{Theorem}[section]
\newtheorem{lemma}[theorem]{Lemma}
\newtheorem{proposition}[theorem]{Proposition}
\newtheorem{corollary}[theorem]{Corollary}
\theoremstyle{definition}
\newtheorem{definition}[theorem]{Definition}
\theoremstyle{remark}
\newtheorem{remark}[theorem]{Remark}
\newcommand{\R}{\mathbb R}
\newcommand{\one}{\mathbf 1}
\newcommand{\tr}{\operatorname{tr}}
\newcommand{\dist}{\operatorname{dist}_1}
\newcommand{\norm}[1]{\left\lVert#1\right\rVert_1}
\newcommand{\abs}[1]{\left|#1\right|}
\newcommand{\pos}[1]{\left(#1\right)_+}
\newcommand{\lean}[1]{\nolinkurl{#1}}
\DeclareMathOperator{\supp}{supp}
\title{Kusner's Conjecture in \(\ell_1^5\)}
\author{Xiao Li
\thanks{Email: \href{mailto:lixiaoms@163.com}{\texttt{lixiaoms@163.com}}.}}
\date{October 8, 2026}

\begin{document}
\maketitle

\begin{abstract}
We prove Kusner's conjecture in $\ell_1^5$: every equilateral subset of real
five-dimensional Manhattan space is finite and has at most ten elements,
and the bound is attained. The argument combines exact energy identities
for ordered coordinate cuts with geometric width bounds on four-point
subsets to exclude eleven equilateral points. A Lean formalization verifies
the theorem for arbitrary real coordinates and every positive common
distance.
\end{abstract}

\section{Introduction and statement}

For $x,y\in\R^d$, write
\[
 \dist(x,y)=\norm{x-y}=\sum_{k=1}^d\abs{x_k-y_k}.
\]
A set is equilateral if there is a real number $\delta>0$ such that
$\dist(x,y)=\delta$ whenever $x$ and $y$ are distinct members. The
equilateral number $e(\ell_1^d)$ is the largest possible cardinality of such
a set. The standard coordinate construction gives $e(\ell_1^d)\geq2d$;
the assertion that equality always holds is usually called Kusner's
conjecture. Equilateral sets in finite-dimensional normed spaces and the
special features of $\ell_p$ spaces are discussed in
\cite{KLS2000,AlonPudlak2003}. Koolen, Laurent and Schrijver established
the conjectured value in dimension four \cite{KLS2000}. The recent
account \cite{GeXuZhou2026} still records the rectilinear equality for
dimensions at most four. We prove the equality in dimension five.

\begin{theorem}\label{thm:main}
Let $S\subseteq\R^5$ and let $\delta>0$. If
\[
 \sum_{k=1}^5\abs{x_k-y_k}=\delta
 \qquad(x,y\in S,\ x\ne y),
\]
then $S$ is finite and $\abs{S}\leq10$. There is a unit-equilateral set
of ten points in $\R^5$. Consequently $e(\ell_1^5)=10$.
\end{theorem}

We prove the exclusion of eleven unit-equilateral points first. All
coordinates in the argument are the five original scalar coordinates.
Auxiliary clamps and projections are introduced explicitly, and their
changes to distance are accounted for. Scalar ties and constant coordinates
are allowed throughout.

The proof has two interacting parts. Ordered scalar cuts give exact trace
and energy identities, which force at least four coordinate ranges to
exceed one half and restrict their extreme pairs to five possible forests.
Four of these possibilities are excluded by explicit prefix profiles,
leaving a matching.

To exclude the remaining matching, we prove that any four original points
have total coordinate width at least $37/24$. This bound uses the same
seven remaining original points, signed overlaps, projections with exactly
accounted displacement, and physical perimeters of plane sections of boxes.
Applied to a suitable four-point subset in the matching, it couples the
five coordinate ranges and forces positive slack away from the otherwise
possible endpoint of the energy inequalities. The resulting contradiction
has an explicit positive rational margin.

\section{Orders, cuts, and exact energy identities}\label{sec:cuts}

Suppose, for a contradiction, that $x_1,\ldots,x_{11}\in\R^5$ satisfy
\begin{equation}\label{eq:unit}
 \dist(x_i,x_j)=1\qquad(i\ne j).
\end{equation}
For coordinate $k$, choose a permutation $p_k$ of the original labels such
that
\[
 x_{p_k(1),k}\leq\cdots\leq x_{p_k(11),k}.
\]
Any consistent order of tied entries may be used. Put
\[
 g_{k,r}=x_{p_k(r+1),k}-x_{p_k(r),k}\geq0,
 \quad A_{k,r}=\{p_k(1),\ldots,p_k(r)\},
 \quad L_k=\sum_{r=1}^{10}g_{k,r}.
\]
Thus $L_k$ is the actual maximum minus minimum of coordinate $k$.

Let $J=\one\one^{\mathsf T}$ and $P=I-J/11$, acting on $\R^{11}$.
Let $D_k$ be the matrix with entries $\abs{x_{i,k}-x_{j,k}}$, and define
\[
 C_k=-\frac12PD_kP.
\]
If $\chi_A$ is the indicator vector of a subset $A$, its cut matrix has
entries $1$ precisely when one label is in $A$ and the other is not. This
matrix is
\[
 \chi_A\one^{\mathsf T}+\one\chi_A^{\mathsf T}
       -2\chi_A\chi_A^{\mathsf T}.
\]
The scalar distance between two ordered entries is the sum of the gaps
whose cuts separate them. Since $P\one=0$, it follows that
\begin{equation}\label{eq:cutmatrix}
 C_k=\sum_{r=1}^{10}g_{k,r}
      (P\chi_{A_{k,r}})(P\chi_{A_{k,r}})^{\mathsf T}\succeq0.
\end{equation}
The complete distance equations give $\sum_kD_k=J-I$ and therefore
\begin{equation}\label{eq:operator}
 \sum_{k=1}^5C_k=\frac12P,
 \qquad 0\preceq C_k\preceq\frac12P,
 \qquad\sum_{k=1}^5\tr C_k=5.
\end{equation}
Writing
\[
 w_r=\frac{r(11-r)}{11},
\]
we also obtain the exact trace formula
\begin{equation}\label{eq:trace}
 T_k:=\tr C_k=\sum_{r=1}^{10}g_{k,r}w_r,
\end{equation}
because $\lVert P\chi_A\rVert_2^2=\abs{A}(11-\abs{A})/11$.

For later use, take any real heights $h_0,\ldots,h_{11}$ with
$h_0=h_{11}=0$ and transport their increments to the original labels:
\[
 v_k(p_k(r))=h_r-h_{r-1}\qquad(1\leq r\leq11).
\]
Then $\sum_iv_k(i)=0$, and its sum on $A_{k,r}$ equals $h_r$. Thus
\begin{equation}\label{eq:profile}
 v_k^{\mathsf T}C_kv_k=\sum_{r=1}^{10}g_{k,r}h_r^2,
 \qquad
 v_k^{\mathsf T}\Bigl(\sum_{\ell=1}^5C_\ell\Bigr)v_k
   =\frac12\sum_{r=1}^{11}(h_r-h_{r-1})^2.
\end{equation}
For a fixed profile put $a(h)=\frac12\sum_r(h_r-h_{r-1})^2$.
Its directed cross energy is
\begin{equation}\label{eq:crossgeneral}
 \Delta(h):=\sum_{k=1}^5\sum_{\ell\ne k}v_k^{\mathsf T}C_\ell v_k
 =5a(h)-\sum_{k=1}^5\sum_{r=1}^{10}g_{k,r}h_r^2.
\end{equation}
The sum is over ordered pairs of coordinates. Each summand is
nonnegative. A constant scalar coordinate still has an order and a
transported vector with the same value of $a(h)$; its own energy is zero.
Hence constant columns are included in these identities without omission.

\section{Triple geometry and the coordinate range bound}\label{sec:range}

\begin{lemma}\label{lem:triple}
For three distinct original points $a,b,c$, their coordinatewise median
$m$ satisfies
$\dist(a,m)=\dist(b,m)=\dist(c,m)=1/2$. Their total coordinate width is
$3/2$.
\end{lemma}

\begin{proof}
In each scalar coordinate, the median lies on all three pairwise intervals.
Consequently
$\dist(a,b)=\dist(a,m)+\dist(b,m)$, and the other two analogous
identities hold. The three left sides are one, so each of the three
distances to $m$ is one half. In one coordinate the sum of the three
pairwise absolute differences is twice its maximum minus minimum.
Summing over coordinates gives total width $3/2$.
\end{proof}

Call a coordinate \emph{wide} if $L_k>1/2$, and put
$t_k=L_k-1/2$ for a wide coordinate. Choose its minimum and maximum
labels $a_k,b_k$. Applying Lemma~\ref{lem:triple} to these two points and
any other label $j$, the median in coordinate $k$ is $x_{j,k}$. Hence
\begin{equation}\label{eq:inner}
 x_{b_k,k}-\frac12\leq x_{j,k}\leq x_{a_k,k}+\frac12
       \qquad(j\notin\{a_k,b_k\}).
\end{equation}
These bounds show that the two extreme labels are unique and that both
outer ordered gaps are at least $t_k$. They also show that the width of
any subset avoiding both extreme labels is at most $1-L_k$.
For every coordinate, wide or not, we therefore have
\begin{equation}\label{eq:outerrange}
 L_k-\frac{g_{k,1}+g_{k,10}}2\leq\frac12.
\end{equation}
Indeed, in a wide column each outer gap is at least $L_k-1/2$;
in a nonwide column the gaps are nonnegative.

We require the following sharper universal range bound.

\begin{proposition}\label{prop:range}
Every original coordinate satisfies $L_k\leq5/6$. Thus every wide excess
satisfies $0<t_k\leq1/3$.
\end{proposition}

\begin{proof}
Only a wide coordinate needs consideration. Translate it so that its
extreme values are $0,L$, and let $a,b$ be their labels. In the remaining
four coordinates, translate and reflect separately so that $a$ is the
zero vector and $b=(\ell_1,\ldots,\ell_4)$, where
\[
 \ell_i\geq0,\qquad \rho:=\sum_{i=1}^4\ell_i=1-L.
\]
For each of the nine other original points, let $u_j$ denote its selected
scalar coordinate. Clamp its remaining coordinates to the box
$\prod_i[0,\ell_i]$, obtaining $y_j$, and write its remainder as $w_j$.
The coordinatewise median of $a,b,x_j$ is $(u_j,y_j)$.
Lemma~\ref{lem:triple} gives
\begin{equation}\label{eq:anchoredmass}
 \norm{w_j}=\frac12,\qquad
 \sum_i y_{j,i}=\frac12-u_j.
\end{equation}

For real numbers define their signed overlap by
\[
 \omega(s,t)=\frac{\abs{s}+\abs{t}-\abs{s-t}}2.
\]
It is the smaller absolute value when $s,t$ have the same sign, and zero
when their signs are opposite. Set $I(v,w)=\sum_i\omega(v_i,w_i)$.
Then
\begin{equation}\label{eq:overlapidentity}
 \norm{v-w}=\norm v+\norm w-2I(v,w).
\end{equation}
Both a scalar clamp and its remainder are monotone, and their increments
sum to the original increment. Thus clamping splits every original
distance exactly into the clamp distance and the remainder distance.
Using \eqref{eq:anchoredmass} for two distinct outside labels yields
\begin{equation}\label{eq:anchoredoverlap}
 I(w_j,w_h)=\frac12\left(
  \norm{y_j-y_h}+\abs{\sum_i(y_{j,i}-y_{h,i})}\right)
 \leq\norm{y_j-y_h}.
\end{equation}
If $w_j,w_h$ have the same nonzero signed ray in coordinate $i$, their
clamps have the same endpoint there. Therefore, for any set $E$ of such
coordinates,
\begin{equation}\label{eq:anchorpaid}
 I(w_j,w_h)+\sum_{i\in E}\ell_i\leq\rho.
\end{equation}

Choose one label attaining the maximum remainder magnitude on each of
the eight signed coordinate rays. There are nine labels, so some label
$j$ was not chosen on any ray. For each nonzero coordinate of $w_j$,
assign its ray's chosen label as owner. Group the coordinates by their
\emph{distinct} owners; call their set $H$, and let $E_h$ be the group
assigned to $h$. We have $h\ne j$, $\abs H\leq4$, and
\[
 \sum_{i\in E_h}\abs{w_{j,i}}\leq I(w_j,w_h).
\]
This uses actual maximum magnitudes on the same rays, so each assigned
coordinate contributes its full magnitude to the overlap. Summing
\eqref{eq:anchorpaid} once for each distinct owner gives
\begin{equation}\label{eq:ownercap}
 \frac12+\sum_{i\in\supp w_j}\ell_i\leq\abs H\rho.
\end{equation}
If $\abs H\leq3$, this implies $1/2\leq3\rho$. If $\abs H=4$,
all four physical coordinates must belong to the support, so the width
sum on the left is $\rho$ and the same conclusion follows. Thus
$1/2\leq3(1-L)$, giving $L\leq5/6$.
Zero box widths cause no difficulty: their remainders and their signed
rays were retained, and their contributions to the subtracted width sum
are simply zero.
\end{proof}

\section{Prefix profiles and wide-coordinate forests}\label{sec:forest}

We now evaluate four particular profiles. Their constants will also be
used in the later matching estimates.

\subsection{Forcing at least four wide coordinates}

First take the symmetric profile whose first five positive heights are
\[
 (h_1,h_2,h_3,h_4,h_5)
   =\frac1{\sqrt{22}}(3,\sqrt{14},\sqrt{26},\sqrt{34},\sqrt{38}),
 \qquad h_{11-r}=h_r.
\]
Its squares equal $w_r-1$ at the eight internal cuts and $w_r-1/2$
at the two outer cuts. Expanding its increment squares gives
\[
 a_M=\frac{204-6\sqrt{14}-4\sqrt{91}-4\sqrt{221}-4\sqrt{323}}{22}.
\]
Equations~\eqref{eq:profile} and \eqref{eq:outerrange} imply
\begin{equation}\label{eq:modifiedcap}
 T_k\leq a_M+L_k-\frac{g_{k,1}+g_{k,10}}2
 \leq\beta_M:=a_M+\frac12<\frac{21}{20}.
\end{equation}
For a nonwide coordinate use instead $h_r=\sqrt{w_r-2/3}$,
with zero endpoint heights. Its first five heights are
$\frac1{\sqrt{33}}(\sqrt8,\sqrt{32},\sqrt{50},\sqrt{62},\sqrt{68})$.
Its half squared increment norm is
\[
 a_N=\frac{260-20\sqrt{31}-4\sqrt{1054}}{33}.
\]
Thus
\begin{equation}\label{eq:nonwidecap}
 T_k\leq a_N+\frac23L_k\leq a_N+\frac13<\frac{10}{11}
          \qquad(L_k\leq1/2).
\end{equation}
All strict numerical inequalities used here and below are certified by
rational square comparisons in Section~\ref{sec:arithmetic}.
If there were at most three wide coordinates, then
\[
 5=\sum_kT_k<3\frac{21}{20}+2\frac{10}{11}
    =\frac{1093}{220}<5,
\]
where replacing a nonwide cap by a wide cap only increases the right
side. This contradiction proves that at least four coordinates are wide.

\subsection{The shifted and unshifted root profiles}

For the remainder of the energy argument use
\[
 h_r=\sqrt{w_r-\frac12},\qquad h_0=h_{11}=0.
\]
The first five heights are
$\frac1{\sqrt{22}}(3,5,\sqrt{37},3\sqrt5,7)$, followed by their
reflection. The squared increments, in positional order, are
\begin{equation}\label{eq:squares}
 s_1,s_2,s_3,s_4,s_5,0,s_5,s_4,s_3,s_2,s_1,
\end{equation}
where
\begin{align}\label{eq:sdefinitions}
 s_1&=\frac9{22},&s_2&=\frac2{11},&
 s_3&=\frac{62-10\sqrt{37}}{22},\notag\\
 s_4&=\frac{82-6\sqrt{185}}{22},&
 s_5&=\frac{94-42\sqrt5}{22}.
\end{align}
They satisfy $s_1>s_2>s_3>s_4>s_5>0$. Put
\begin{equation}\label{eq:alpha}
 \alpha=\sum_{i=1}^5s_i
   =\frac{251-10\sqrt{37}-6\sqrt{185}-42\sqrt5}{22}<\frac23.
\end{equation}
The own-coordinate energy is $T_k-L_k/2$. Writing
$W=\sum_kL_k$, the exact directed cross energy is consequently
\begin{equation}\label{eq:delta}
 \Delta=5\alpha+\frac W2-5\geq0.
\end{equation}
Also $T_k\leq\alpha+L_k/2$, so $W\geq10-10\alpha>10/3$.
If
\[
 H=\sum_{k\text{ wide}}t_k,
\]
then $W\leq5/2+H$, whence
\begin{equation}\label{eq:Hlower}
 H>\frac56.
\end{equation}

The unshifted root profile $h_r=\sqrt{w_r}$ has half squared increment
norm
\begin{equation}\label{eq:alpha0}
 \alpha_0=
 \frac{190-12\sqrt5-24\sqrt3-8\sqrt{42}-4\sqrt{210}}{11}
 <\frac{43}{40}.
\end{equation}
Its own energy is $T_k$, and its outer increment squares are $10/11$.
Its total directed cross energy is exactly $5\alpha_0-5$.

\subsection{The forest and its isolated edges}

Associate to each wide coordinate the edge $\{a_k,b_k\}$ on the eleven
original labels. No two coordinates have the same extreme pair: their
two ranges alone would give distance greater than one for that pair.
There is no cycle either. If a cycle uses $m$ wide edges, the variation
of each associated scalar coordinate along the entire cycle is at least
twice its range. Summing those $m$ coordinate variations gives more than
$m$, although the cycle has $m$ original unit-distance edges. This is
impossible. Thus the wide-edge graph is a simple forest.

For every zero-sum vector $v$, its energy in a wide target coordinate is
bounded below by its two outer cuts:
\begin{equation}\label{eq:outerenergy}
 v^{\mathsf T}C_\ell v\geq
 t_\ell\bigl(v(a_\ell)^2+v(b_\ell)^2\bigr).
\end{equation}
Here the last prefix has sum $-v(b_\ell)$, explaining why both endpoint
squares occur with the same factor $t_\ell$.
Apply this to the unshifted root vectors. An endpoint common to the
source row and target edge has square $10/11$. If $\deg(a)$ denotes
degree in the wide forest, summing over the ordered source-target pairs
gives
\begin{equation}\label{eq:rootledger}
 5\alpha_0-5\geq\frac{10}{11}
 \sum_{\ell\text{ wide}}t_\ell
  \bigl(\deg(a_\ell)+\deg(b_\ell)-2\bigr).
\end{equation}
For a nonisolated edge the degree factor is at least one; for an
isolated edge it is zero. Hence the excess sum $H_{\rm non}$ of the
nonisolated edges satisfies
\[
 H_{\rm non}\leq\frac{11}{10}(5\alpha_0-5)<\frac{33}{80}.
\]
Together with \eqref{eq:Hlower}, this gives
\begin{equation}\label{eq:isolatedmass}
 H_{\rm iso}>\frac56-\frac{33}{80}=\frac{101}{240}.
\end{equation}
Each individual excess is at most $1/3=80/240$ by
Proposition~\ref{prop:range}; thus there are at least two isolated edges.

There are four or five wide edges. A nonisolated edge has a nonisolated
neighbor, so there cannot be exactly one nonisolated edge. The complete
list of possibilities is therefore
\begin{equation}\label{eq:cases}
 (m,i)=(4,2),(4,4),(5,2),(5,3),(5,5),
\end{equation}
where $m$ is the number of wide edges and $i$ the number of isolated
ones. In $(4,2)$ and $(5,3)$ the nonisolated part is a two-edge path.
In $(5,2)$ it is a three-edge tree: it must be connected, since two
nontrivial components would require at least four nonisolated edges;
it then has four vertices because it is a forest. The cases $(4,4)$
and $(5,5)$ are matchings.

\section{Excluding four of the five forest shapes}\label{sec:exclusions}

We use the shifted profile of \eqref{eq:squares}. Its actual extreme
entries have square $s_1=9/22$. The same counting as in
\eqref{eq:rootledger} consequently yields
\begin{equation}\label{eq:shiftednonisolated}
 \Delta\geq\frac9{22}H_{\rm non}.
\end{equation}
No assertion about entries at other labels is needed for this inequality.

\subsection{Two isolated edges}

Suppose first that $(m,i)=(4,2)$. The two isolated excesses sum to at most
$2/3$, so $H_{\rm non}\geq H-2/3\geq W-19/6$. The four wide ranges
and the nonwide range satisfy $W\leq4(5/6)+1/2=23/6$. Therefore
\[
 \Delta\geq\frac9{22}\left(W-\frac{19}{6}\right),
 \qquad
 \Delta<\frac W2-\frac53
\]
by \eqref{eq:alpha} and \eqref{eq:delta}. The lower bound minus the
strict upper bound is
\[
 \frac{49}{132}-\frac W{11}\geq\frac1{44}>0,
\]
a contradiction.

Now suppose $(m,i)=(5,2)$. The three-edge tree has four vertices and the
two isolated edges have four more, leaving three original labels that
are extreme in no coordinate. Call them core labels. Their total
coordinate width is $3/2$ by Lemma~\ref{lem:triple}. In each coordinate
\eqref{eq:inner} bounds their width by $1-L_k$, so
\[
 \frac32\leq5-W,\qquad W\leq\frac72,
 \qquad H=W-\frac52\leq1.
\]
Again the isolated excesses sum to at most $2/3$. Thus
\[
 \Delta\geq\frac9{22}\left(H-\frac23\right),
 \qquad
 \Delta<\frac H2-\frac5{12}.
\]
Their difference is $19/132-H/11\geq7/132>0$. This excludes $(5,2)$.

\subsection{A finite positional estimate}\label{subsec:positions}

Let $q(r)$ be the squared increment at position $r$ in
\eqref{eq:squares}. The smallest $2,4,6,8$ entries of its eleven-entry
multiset have respective sums
\begin{equation}\label{eq:minimalsums}
 M_2=s_5,\quad M_4=2s_5+s_4,\quad
 M_6=2s_5+2s_4+s_3,\quad
 M_8=2s_5+2s_4+2s_3+s_2.
\end{equation}
For any $2r$ distinct original labels their squared entries in a given
row sum to at least $M_{2r}$. Positions remain distinct even when their
scalar values are tied.

Define
\begin{equation}\label{eq:kappas}
 \kappa_1=s_5,\quad \kappa_2=s_5+s_4,\quad
 \kappa_3=s_4+s_3,\quad \kappa_4=s_3+s_2.
\end{equation}
Suppose $u_1\geq\cdots\geq u_r\geq0$, where $r=3$ or $4$, and
each $u_j$ is assigned to two distinct labels, all $2r$ labels being
distinct. If $Q_j$ is the sum of their two squared entries, then
\begin{equation}\label{eq:pairedrearrangement}
 \sum_{j=1}^ru_jQ_j\geq\sum_{j=1}^r\kappa_j u_j.
\end{equation}
Indeed, write the left side as
\[
 u_r\sum_{j=1}^rQ_j+
 \sum_{a=1}^{r-1}(u_a-u_{a+1})\sum_{j=1}^aQ_j.
\]
Every partial sum contains $2a$ distinct positions, so it is at least
$M_{2a}$. Substitution of \eqref{eq:minimalsums} proves
\eqref{eq:pairedrearrangement}. This argument is a lower relaxation of
actual indexed positions; it never chooses new scalar orders or assumes
that independently favorable placements can be realized.

We shall also need a weighted threshold version. For distinct positions
$f(1),\ldots,f(n)$, nonnegative weights $u_j\leq M$, and $a\geq0$,
\begin{equation}\label{eq:threshold}
 \sum_ju_jq(f(j))\geq
 a\sum_ju_j-M\sum_{r=1}^{11}\pos{a-q(r)}.
\end{equation}
To prove it, write
\[
 u_jq(f(j))=au_j-u_j(a-q(f(j))).
\]
Bound the subtracted term by $M\pos{a-q(f(j))}$, and use distinctness
to bound the sum of these nonnegative discounts by the full positional
sum. If one specified weight equals $S\geq M$ and all other weights are
at most $M$, the same proof, separating the extra weight $S-M$, gives
\begin{equation}\label{eq:onelarge}
 \sum_ju_jq(f(j))\geq
 a\sum_ju_j-M\sum_{r=1}^{11}\pos{a-q(r)}-(S-M)a.
\end{equation}
The last term is valid since $q(r)\geq0$, so every discount is at most
$a$. Directly from \eqref{eq:squares},
\begin{equation}\label{eq:discounts}
 \sum_r\pos{s_3-q(r)}=5s_3-2s_4-2s_5=:D_3,
 \quad
 \sum_r\pos{s_4-q(r)}=3s_4-2s_5=:D_4.
\end{equation}

\subsection{The matching of four wide edges}\label{subsec:fourmatching}

Suppose $(m,i)=(4,4)$, and order its wide excesses as
$t_1\geq t_2\geq t_3\geq t_4>0$. Write $L_0\leq1/2$ for the
remaining range. The eight extreme labels leave three core labels.
Their width in a wide coordinate is at most $1-L_k$ and in the other
coordinate at most $L_0$. Lemma~\ref{lem:triple} therefore gives
\[
 \frac32\leq\sum_{k=1}^4(1-L_k)+L_0
      =2-H+L_0,
 \qquad H:=\sum_{k=1}^4t_k\leq1.
\]

In each wide source row the three other edges contribute their two
distinct endpoint positions through \eqref{eq:outerenergy}. Applying
\eqref{eq:pairedrearrangement} with $r=3$ gives the three coefficients
$\kappa_1,\kappa_2,\kappa_3$ in descending weight order after omitting
that row's own edge. In the nonwide source row all four edges contribute,
and the coefficients are $\kappa_1,\ldots,\kappa_4$. Summing these
five actual rows gives
\begin{equation}\label{eq:Bdefinitions}
 \begin{split}
 B_1&=4\kappa_1,\qquad B_2=\kappa_1+3\kappa_2,\qquad
 B_3=2\kappa_2+2\kappa_3,\\
 B_4&=3\kappa_3+\kappa_4,\qquad B_5=4\kappa_4,
 \end{split}
\end{equation}
and
\begin{equation}\label{eq:fourmatchinglower}
 \Delta\geq\sum_{j=1}^4B_jt_j.
\end{equation}
For example, $t_2$ receives $\kappa_1$ from the first wide row,
$\kappa_2$ from the third and fourth wide rows, and another $\kappa_2$
from the nonwide row. This explains the ordered-row count and its factor
three. Every edge weight has already been attached to its two original
endpoint labels in \eqref{eq:pairedrearrangement}.

The coefficients satisfy
\begin{equation}\label{eq:Bsigns}
 B_1<B_2<B_3<B_4<B_5,\qquad B_3<\frac{17}{64},
 \qquad B_4<\frac12,
 \qquad B_5+B_3>1.
\end{equation}
Their exact certificates are proved in Section~\ref{sec:arithmetic}.
Since $t_1,t_2\leq1/3$ and $t_4\geq0$, increasing coefficients give
\begin{equation}\label{eq:fourcapped}
 \sum_{j=1}^4B_jt_j
 \geq\frac{B_1+B_2}{3}+B_3\left(H-\frac23\right).
\end{equation}
For clarity, the difference between the two sides is
\[
 (B_3-B_1)(1/3-t_1)+(B_3-B_2)(1/3-t_2)
      +(B_4-B_3)t_4\geq0.
\]
On the other hand $W=2+H+L_0\leq5/2+H$, so
\begin{equation}\label{eq:matchingupper}
 \Delta\leq5\alpha-\frac{15}{4}+\frac H2.
\end{equation}
Because $B_3<1/2$ and $H\leq1$, the difference between the lower bound
in \eqref{eq:fourcapped} and the upper bound in
\eqref{eq:matchingupper} is at least
\begin{equation}\label{eq:G1}
 G_1:=\frac{B_1+B_2+B_3}{3}-5\alpha+\frac{13}{4}
 =\frac{-3825+260\sqrt{37}+96\sqrt{185}+420\sqrt5}{132}
 >\frac{17549}{1650000}>0.
\end{equation}
This excludes the four-edge matching.

\subsection{Three isolated edges and a two-edge path}\label{subsec:path}

It remains in this section to exclude $(m,i)=(5,3)$. Let the three
isolated excesses be $u_1,u_2,u_3$, with sum $U\leq1$. Write the
path as $A$--$B$--$C$, with excesses $a,c$ and sum $S=a+c\leq2/3$.
There are nine distinct endpoint labels: six isolated endpoints and
the three path vertices. Each wide source row has seven other endpoint
labels in distinct internal positions.

Consider the source row belonging to edge $AB$. In the target edge
$BC$, the shared label $B$ is one of the source row's own extremes,
so its square is $s_1$. Its contribution is $cs_1$. The other seven
labels have weights $u_1,u_1,u_2,u_2,u_3,u_3,c$, each at most $1/3$.
Equation~\eqref{eq:threshold} with threshold $s_3$ bounds their
contribution below by $s_3(2U+c)-D_3/3$. The analogous statement for
row $BC$ has $as_1$ and $s_3(2U+a)-D_3/3$. The two path rows
together therefore contribute at least
\begin{equation}\label{eq:pathrows}
 4s_3U+(s_1+s_3)S-\frac23D_3.
\end{equation}

In the row of an isolated edge with excess $u$, the seven other labels
have total weight $2(U-u)+2S$. The shared path vertex $B$ occurs
\emph{once}, with weight $S=a+c$; the weights at $A,C$ are $a,c$.
The other four weights are the endpoints of the two other isolated
edges. If $S\leq1/3$, all seven weights are at most $1/3$, and the
threshold-$s_4$ discount is at most $D_4/3$.
If $S\geq1/3$, only the single shared weight can exceed $1/3$.
Equation~\eqref{eq:onelarge} bounds the discount by
\[
 \frac{D_4}{3}+(S-1/3)s_4
    =s_4S+\frac23(s_4-s_5).
\]
This bound holds wherever the shared label is actually positioned.

Summing these three isolated rows with \eqref{eq:pathrows} gives the
following lower bounds for the same original $\Delta$:
\begin{align}\label{eq:pathsmall}
 F_{\rm small}(U,S)
 &=4(s_3+s_4)U+(s_1+s_3+6s_4)S
       -\frac{5s_4+10s_3-10s_5}{3}
       &&(S\leq1/3),\\
 F_{\rm large}(U,S)
 &=4(s_3+s_4)U+(s_1+s_3+3s_4)S
       -\frac{2s_4+10s_3-10s_5}{3}
       &&(S\geq1/3).\label{eq:pathlarge}
\end{align}
The bounds agree at $S=1/3$. Since all five coordinates are wide,
\begin{equation}\label{eq:pathexact}
 \Delta=5\alpha-\frac{15}{4}+\frac{U+S}{2}.
\end{equation}
The coefficient signs needed to compare these affine functions are
\begin{equation}\label{eq:pathsigns}
 4(s_3+s_4)-\frac12<0,
 \qquad s_1+s_3+3s_4-\frac12>0.
\end{equation}
In the small-$S$ case, replacing $U\leq1$ by $1$ and $S\geq0$ by
zero can only decrease $F_{\rm small}-\Delta$. Its resulting value is
exactly $G_1$. In the large-$S$ case, replacing $U$ by $1$ and $S$ by
$1/3$ can only decrease $F_{\rm large}-\Delta$. Its value is
\[
 G_1+\frac13\left(s_1+s_3+6s_4-\frac12\right)>G_1.
\]
Both are positive by \eqref{eq:G1} and \eqref{eq:pathsigns},
contradicting $\Delta\geq F$. This excludes $(5,3)$.

Thus the only possible wide forest consists of five disjoint edges.
Before excluding that matching, we establish a bound for every four-point
subset using all seven of its remaining original labels.

\section{Four original points and their seven remaining labels}\label{sec:quartet}

Fix any four distinct original labels, denoted $q_0,q_1,q_2,q_3$.
In each physical coordinate sort their four values. Its first and last
gaps are singleton cuts, assigned respectively to the least and greatest
label, and its middle gap is one of the three balanced two--two cuts.
Let $b_i$ be the total singleton width assigned to $q_i$, and let
$c_1,c_2,c_3$ be the total widths of the three balanced cuts.

For complementary pairs, the balanced cuts make equal contributions.
Subtracting their original unit-distance equations gives
\[
 b_0+b_1=b_2+b_3,\quad b_0+b_2=b_1+b_3,
 \quad b_0+b_3=b_1+b_2.
\]
Thus $b_0=b_1=b_2=b_3=:b\geq0$. Each of the six pairs crosses
two of the three balanced cuts. Their equations now imply
\begin{equation}\label{eq:quartetcut}
 c_1=c_2=c_3=:c=\frac12-b\geq0.
\end{equation}
The actual total coordinate width $W_Q$ of the four points is therefore
\begin{equation}\label{eq:quartetwidth}
 W_Q=4b+3c=\frac32+b.
\end{equation}
These identities also hold when a sorted gap is zero.

\subsection{Clamps, signed owners, and all four cross equations}

In each coordinate clamp every original point to the interval between
the second and third selected values and subtract that interval's
midpoint. Let its halfwidth be $\lambda_k\geq0$. Reflect the coordinate
if necessary and assign its balanced type $\tau(k)\in\{1,2,3\}$ so
that the selected clamps are $H_{i,\tau(k)}\lambda_k$, where
\begin{equation}\label{eq:Hmatrix}
 H=\begin{pmatrix}
 1&1&1\\1&-1&-1\\-1&1&-1\\-1&-1&1
 \end{pmatrix},\qquad
 H^{\mathsf T}\one=0,\quad H^{\mathsf T}H=4I_3.
\end{equation}
Every balanced partition of four labels is a column of this matrix,
up to sign, so this uses only actual coordinate reflections. Coordinates
with zero middle gap are still assigned a type. Write the resulting
decomposition of each point as
\[
 x_j=y_j+w_j+\text{midpoint},
 \qquad -\lambda_k\leq y_{j,k}\leq\lambda_k,
\]
with the understood coordinate reflections on the left side. They
preserve the physical $\ell_1$ distance. As before, monotonicity of the
clamp and its remainder gives the exact splitting
\begin{equation}\label{eq:quartetsplit}
 \dist(x_j,x_\ell)=\norm{y_j-y_\ell}+\norm{w_j-w_\ell}
 \quad\text{for every original pair.}
\end{equation}
In particular a zero halfwidth forces $y_{j,k}=0$ but does not force
$w_{j,k}=0$.

Let $u_i=w_{q_i}$ be the four selected tails. Their masses are
$\norm{u_i}=b$. In one coordinate only the selected least value can
have a negative tail and only the greatest can have a positive tail.
Thus two owners never occupy the same nonzero signed ray; opposite
rays of one coordinate are allowed. Also
\begin{equation}\label{eq:lambdatotal}
 \sum_{k:\tau(k)=t}\lambda_k=\frac c2,
 \qquad
 \Lambda:=\sum_k\lambda_k=\frac{3c}{2}
      =\frac34-\frac{3b}{2}.
\end{equation}
The norm of each selected point relative to these midpoints is
$b+\Lambda=3/4-b/2$. This quantity differs from $\Lambda$.

There are exactly seven other original labels. For each such label $j$
define
\[
 A_{i,j}=I(u_i,w_j),\quad J_j=\sum_{i=0}^3A_{i,j},\quad
 r_j=\norm{w_j},\quad
 Q_{t,j}=\sum_{k:\tau(k)=t}y_{j,k}.
\]
The signed overlap of \eqref{eq:overlapidentity} satisfies
\begin{equation}\label{eq:agreementbounds}
 0\leq A_{i,j}\leq b.
\end{equation}
In a selected clamp coordinate, the distance to an arbitrary clamp is
$\lambda_k-H_{i,\tau(k)}y_{j,k}$. The original pair $(q_i,j)$
therefore gives, for each of the four values of $i$,
\begin{equation}\label{eq:crossquartet}
 1=b+\Lambda+r_j-(HQ_j)_i-2A_{i,j}.
\end{equation}
Summing these four equations, and then multiplying them by
$H^{\mathsf T}$, yields respectively
\begin{equation}\label{eq:outsiderequations}
 r_j=\delta+\frac{J_j}{2},\qquad
 Q_j=-\frac12H^{\mathsf T}A_j,
 \qquad \delta=\frac14+\frac b2.
\end{equation}
Here $A_j$ is the four-entry agreement vector. Thus no equality of the
original outsider norms has been assumed. Equations
\eqref{eq:outsiderequations} are consequences of all four original
cross distances.

\subsection{Co-occurrence on physical signed rays}\label{subsec:cooccurrence}

For two remaining labels $j,\ell$ the disjoint signed owner rays imply
\begin{equation}\label{eq:cooccurrence}
 I(w_j,w_\ell)\geq
    \sum_{i=0}^3\pos{A_{i,j}+A_{i,\ell}-b}.
\end{equation}
Here is a coordinatewise proof that also accounts for opposite owner
rays in a single coordinate. For nonnegative $u,a,a'$,
\[
 \min(u,a,a')\geq\min(u,a)+\min(u,a')-u.
\]
Equivalently, in signed-overlap notation, for any real $u,a,a'$,
\[
 \min\{\omega(u,a),\omega(u,a')\}
 \geq\omega(u,a)+\omega(u,a')-\abs u.
\]
This follows because each overlap is at most $\abs u$. Sum over
coordinates for one owner $u_i$, and use nonnegativity of the left
side to bound it below by
$\pos{A_{i,j}+A_{i,\ell}-\norm{u_i}}$.
Now sum over owners. If $w_j,w_\ell$ have opposite signs in a
coordinate, all the minima in that coordinate are zero. If they have
the same sign, only the one possible owner on that signed ray can
contribute, and its minimum overlap is at most
$\omega(w_{j,k},w_{\ell,k})$. Thus no physical overlap is counted
twice. This proves \eqref{eq:cooccurrence} using $\norm{u_i}=b$.

\section{Projection with exactly accounted displacement}\label{sec:projection}

\begin{lemma}\label{lem:boxprojection}
Let $\lambda_1,\ldots,\lambda_n\geq0$ and
$v\in\prod_i[-\lambda_i,\lambda_i]$. For every feasible sum
$-\sum_i\lambda_i\leq T\leq\sum_i\lambda_i$ there exists a vector
$z$ in the same box such that
\[
 \sum_i z_i=T,\qquad
 \sum_i\abs{z_i-v_i}=\abs{\sum_i v_i-T}.
\]
In particular, the distance to the zero-sum section is exactly
$\abs{\sum_i v_i}$.
\end{lemma}

\begin{proof}
Write $Q=\sum_i v_i$ and $L=\sum_i\lambda_i$. If $Q=T$, take
$z=v$. If $T<Q$, then $Q+L>0$ and
\[
 \theta=\frac{Q-T}{Q+L}\in[0,1],\qquad
 z_i=(1-\theta)v_i-\theta\lambda_i
\]
defines a point in the box with sum $T$. Each $z_i\leq v_i$, so the
sum of absolute changes is $Q-T$. The case $Q<T$ follows by negating
$v,T,z$. Finally every vector with sum $T$ is at distance at least
$\abs{Q-T}$, by the triangle inequality for the sum of coordinate
differences. This proves exactness, including zero widths.
\end{proof}

\begin{lemma}\label{lem:faceprojection}
Assume in addition that $n\geq2$ and all $\lambda_i>0$. If $v$ touches
a box face, a closest projection to the zero-sum section can be chosen
on that section's relative boundary, with the same displacement
$\abs{\sum_i v_i}$.
\end{lemma}

\begin{proof}
After negation we may suppose $v_i=\lambda_i$ for some $i$. Put
$M=\sum_{h\ne i}\lambda_h$. If $\lambda_i\leq M$, keep
$z_i=\lambda_i$ and use Lemma~\ref{lem:boxprojection} on the other
coordinates with target sum $-\lambda_i$. Its displacement is exactly
$\abs{\sum_h v_h}$, and $z$ still touches the $i$th face.
If $\lambda_i>M$, take instead
\[
 z_i=M,\qquad z_h=-\lambda_h\quad(h\ne i).
\]
This vector belongs to the box, has zero sum, and is coordinatewise at
most $v$. Also $\sum_hv_h\geq\lambda_i-M>0$, so its displacement is
again $\abs{\sum_hv_h}$. It touches every other face. Such another
coordinate exists and has positive width. A saturated coordinate is a
true relative-boundary condition here: the zero vector is strictly
inside all the box inequalities, and each coordinate functional is
nonconstant on the zero-sum subspace when $n\geq2$.
\end{proof}

Apply Lemma~\ref{lem:boxprojection} separately to the coordinates of
each of the three actual balanced types. For each remaining original
label $j$ obtain $z_j$ in the physical section
\begin{equation}\label{eq:section}
 K=\left\{z\in\R^5:
   \abs{z_k}\leq\lambda_k\ (1\leq k\leq5),\quad
   \sum_{k:\tau(k)=t}z_k=0\ (t=1,2,3)\right\}.
\end{equation}
Its displacement is exactly
\begin{equation}\label{eq:payment}
 E_j:=\norm{z_j-y_j}=\sum_{t=1}^3\abs{Q_{t,j}}
   =\frac12\sum_{t=1}^3\abs{(H^{\mathsf T}A_j)_t}.
\end{equation}
Zero-width coordinates of $z_j,y_j$ are zero; their original tails and
agreements remain in \eqref{eq:outsiderequations} and
\eqref{eq:cooccurrence}.

\subsection{The eight agreement forms and the pair table}

For a real four-vector $A$ let $J(A)=\sum_iA_i$, and put
\[
 E(A)=\frac12\left(
 \abs{A_0+A_1-A_2-A_3}
 +\abs{A_0-A_1+A_2-A_3}
 +\abs{A_0-A_1-A_2+A_3}\right).
\]
The following exact identity is useful:
\begin{equation}\label{eq:formidentity}
 F(A):=\frac{J(A)}2+E(A)
   =\max_{0\leq i\leq3}\{2A_i,\ J(A)-2A_i\}.
\end{equation}
To verify it, sort the four entries as $a\leq b'\leq c'\leq d$.
The three absolute expressions are indexed by the three unordered
balanced partitions and are therefore unchanged as a multiset by
permuting the entries. Their half sum becomes
\[
 d-a+\frac12\abs{a+d-b'-c'}.
\]
Adding half the total gives
$\max(2d,a+b'+c'+d-2a)$, precisely the right side of
\eqref{eq:formidentity}. This proof includes all ties.

Call $2A_i$ a plus form at index $i$ and $J(A)-2A_i$ a minus form at
index $i$. Choose one maximizing form at each of the seven labels,
once and for all. The coefficient vector of a plus form is $2$ at its
index and zero elsewhere. The coefficient vector of a minus form is
$-1$ at its index and $1$ elsewhere.

For $0\leq u,v\leq b$ and real coefficients $a,d$, the function
\[
 au+dv-2\pos{u+v-b}
\]
is linear on each triangle into which $u+v=b$ divides the square.
Checking their vertices gives the upper bound
\begin{equation}\label{eq:slotcap}
 au+dv-2\pos{u+v-b}
   \leq b\max(0,a,d,a+d-2).
\end{equation}
For $b=0$ the assertion is immediate. For $b>0$, divide the variables
by $b$; the four vertex values are $0,a,d,a+d-2$, which proves it.
Apply this independently to the four agreement slots and sum. For any
two chosen forms the resulting total cap is $b$ times the following
number:
\begin{center}
\begin{tabular}{lcc}
\toprule
Chosen forms & Equal indices & Different indices\\
\midrule
plus / plus & $2$ & $4$\\
minus / minus & $3$ & $4$\\
plus / minus & $5$ & $4$\\
\bottomrule
\end{tabular}
\end{center}
For example, equal-index plus/minus forms have cap $2$ at their common
index and cap $1$ at each of the other three indices, totaling $5$.
Different-index plus/minus forms have cap $2$ at the plus index, zero
at the minus index, and $1$ at each remaining index, totaling $4$.
For equal-index minus/minus forms the common-index cap is zero and
the other three caps are $1$. The other entries follow from the same
four-vertex rule. Thus the table bounds actual forms, without any
simultaneous-attainment assertion.

From \eqref{eq:quartetsplit}, \eqref{eq:outsiderequations}, the two
projection displacements, and \eqref{eq:cooccurrence}, we obtain
\begin{align}\label{eq:pairprojected}
 \norm{z_j-z_\ell}
 &\geq1-2\delta-
 \left(F(A_j)+F(A_\ell)
       -2\sum_i\pos{A_{i,j}+A_{i,\ell}-b}\right)\notag\\
 &\geq\frac12-b-b\,C(j,\ell),
\end{align}
where $C(j,\ell)$ is the corresponding table entry. In particular,
\begin{equation}\label{eq:pairseparation}
 \norm{z_j-z_\ell}\geq\frac12-6b\qquad(j\ne\ell).
\end{equation}
To see the first line directly, the original distance equation gives
$\norm{y_j-y_\ell}=1-r_j-r_\ell+2I(w_j,w_\ell)$.
Subtract $E_j+E_\ell$ by the triangle inequality and substitute
$r_j=\delta+J_j/2$ and $E_j=E(A_j)$.

\subsection{A seven-label cyclic lower bound}\label{subsec:cyclelower}

In any cycle on the same seven labels, there is an edge joining forms
of the same sign, since signs cannot alternate around an odd cycle.
If its indices are also equal, its table cap is at most $3$; the other
six caps are at most $5$, giving total at most $33$.
If its indices differ, its cap is $4$. There must be a second edge
whose indices differ: otherwise equality of indices along the other
six edges would force equality across the first edge too. The second
cap is also $4$, while the remaining five are at most $5$.
Again the total is at most $33$. Summing \eqref{eq:pairprojected} proves
\begin{equation}\label{eq:cyclelower}
 \sum_{j\text{ in any cyclic order}}\norm{z_j-z_{j+1}}
     \geq\frac72-40b.
\end{equation}
The forms are chosen per label, so the same choices work in whatever
cyclic order is later supplied by geometry.

\section{Physical boundary tours and the quartet width bound}\label{sec:perimeter}

We now prove
\begin{proposition}\label{prop:quartetbound}
Every four-point subset of the eleven original points has
$b\geq1/24$ and total coordinate width $W_Q\geq37/24$.
\end{proposition}

Assume for a contradiction that $b<1/24$. Then $c=1/2-b>0$, and
each balanced type has positive total halfwidth $c/2$. Let $n_t$ be
the number of its positive physical halfwidths. The type supports are
disjoint and there are only five coordinates, so, after permuting the
three types, the possibilities for $(n_1,n_2,n_3)$ are
\begin{equation}\label{eq:allocations}
 (1,1,1),\quad(2,1,1),\quad(3,1,1),\quad(2,2,1).
\end{equation}
Every zero-width slot has $z_k=0$ in the auxiliary section, irrespective
of its original tail. The pair lower bound \eqref{eq:pairseparation}
now gives
\begin{equation}\label{eq:quarterseparation}
 \norm{z_j-z_\ell}>\frac14\qquad(j\ne\ell).
\end{equation}

For $(1,1,1)$ the section $K$ is a point, contradicting
\eqref{eq:quarterseparation}. For $(2,1,1)$ its only nonzero
coordinates have the form $(r,-r)$, with
$\abs r\leq A:=\min(\lambda_1,\lambda_2)$. Their physical distance
is $2\abs{r-r'}$ and the segment diameter is $4A\leq c\leq1/2$,
since $\lambda_1+\lambda_2=c/2$. Three ordered points on a segment
with successive distances greater than $1/4$ would have endpoint
distance greater than $1/2$. Three of the seven labels already give
a contradiction. No boundary condition was used in these two cases.

\subsection{A genuine positive face for every remaining label}

In each of the last two allocations in \eqref{eq:allocations}, all
five physical coordinates have positive halfwidth. This fact must be
established before passing from a tail to an intrinsic boundary.
If a balanced type has a single coordinate, its halfwidth is $c/2$.
Equation~\eqref{eq:outsiderequations}, together with
\eqref{eq:agreementbounds}, gives
\[
 \abs{Q_{t,j}}\leq b<\frac c2.
\]
Indeed each column of $H$ has two plus and two minus entries, so the
absolute difference of the corresponding two agreement sums is at
most $2b$. Thus a singleton type cannot carry a nonzero original tail:
such a tail would force its clamp to an endpoint, making
$\abs{Q_{t,j}}=c/2$.

On the other hand \eqref{eq:outsiderequations} gives
$\norm{w_j}=\delta+J_j/2>0$. Hence some actual positive-width
coordinate of a type with at least two coordinates has a nonzero tail.
Its clamp touches a face. Apply Lemma~\ref{lem:faceprojection} on
the actual positive coordinates of that type, and use
Lemma~\ref{lem:boxprojection} on the other types. The resulting
$z_j\in K$ still has exactly the displacement \eqref{eq:payment},
and it touches a positive face of a nonsingleton section.
The face witness is therefore inside the active type, rather than an
outside zero coordinate. We use these choices for the seven labels.

\subsection{The two paired coordinates in each of two types}

For allocation $(2,2,1)$ the section, after an actual coordinate
permutation, consists of
\[
 (r,-r,s,-s,0),\qquad \abs r\leq A,\quad\abs s\leq B,
\]
where each of $A,B$ is the smaller of its type's two halfwidths.
The physical metric on this rectangle is
\begin{equation}\label{eq:pairedmetric}
 2\abs{r-r'}+2\abs{s-s'}.
\end{equation}
Each selected $z_j$ lies on a true side, by the positive-face conclusion
above. Traversing the four literal vertices
$(A,-A,B,-B,0)$, $(-A,A,B,-B,0)$,
$(-A,A,-B,B,0)$, $(A,-A,-B,B,0)$ gives side lengths
$4A,4B,4A,4B$. Since $A,B\leq c/4$, the complete tour length is
\begin{equation}\label{eq:doubleperimeter}
 8(A+B)\leq4c=2-4b.
\end{equation}
The factors two in \eqref{eq:pairedmetric} come from two separate
physical coordinates in each pair.

\subsection{The three-coordinate plane section}

For allocation $(3,1,1)$ the two singleton auxiliary coordinates are
zero. The remaining section is the literal plane polygon
\begin{equation}\label{eq:planesection}
 z_1+z_2+z_3=0,\qquad
 \abs{z_1}\leq a,\quad\abs{z_2}\leq d,\quad\abs{z_3}\leq f,
 \qquad a+d+f=\frac c2.
\end{equation}
We give its physical boundary tour explicitly. On an edge with one
coordinate fixed, the other two coordinate differences are opposites.
Its physical distance is therefore \emph{twice} the absolute change
of either free coordinate.

If all three widths satisfy the triangle inequalities, use the six
vertices, in order,
\begin{equation}\label{eq:hexvertices}
 \begin{aligned}
 v_0&=(a,f-a,-f),& v_1&=(a,-d,d-a),\\
 v_2&=(d-f,-d,f),& v_3&=(-a,a-f,f),\\
 v_4&=(-a,d,a-d),& v_5&=(f-d,d,-f).
 \end{aligned}
\end{equation}
They satisfy all the box inequalities. Their six edge lengths are
\begin{equation}\label{eq:hexlengths}
 2(d+f-a),\ 2(a+f-d),\ 2(a+d-f),\
 2(d+f-a),\ 2(a+f-d),\ 2(a+d-f).
\end{equation}
Consequently the total physical length is $4(a+d+f)=2c$.
Every saturated box face is exactly its corresponding segment in
this list. For example, on $z_1=a$ the plane and box inequalities
give $-d\leq z_2\leq f-a$, the segment from $v_1$ to $v_0$.
For completeness the six face intervals are
\begin{center}
\begin{tabular}{ccc}
\toprule
Fixed coordinate & Free coordinate & Its interval\\
\midrule
$z_1=a$ & $z_2$ & $[-d,f-a]$\\
$z_2=-d$ & $z_1$ & $[d-f,a]$\\
$z_3=f$ & $z_1$ & $[-a,d-f]$\\
$z_1=-a$ & $z_2$ & $[a-f,d]$\\
$z_2=d$ & $z_1$ & $[-a,f-d]$\\
$z_3=-f$ & $z_1$ & $[f-d,a]$\\
\bottomrule
\end{tabular}
\end{center}
The remaining coordinate is fixed by the plane equation. This verifies
coverage of all true faces, including endpoint coincidences. If a
triangle inequality is an equality, the corresponding edges in
\eqref{eq:hexlengths} have zero length and are retained as constant
segments.

If a width dominates, permute the actual coordinates and suppose
$a>d+f$. Then the plane implies $\abs{z_1}\leq d+f<a$; the
first coordinate's original faces do not meet the section. The actual
section is the rectangle with vertices
\begin{equation}\label{eq:dominantvertices}
 (d+f,-d,-f),\ (f-d,d,-f),\
 (-d-f,d,f),\ (d-f,-d,f).
\end{equation}
The four physical lengths are $4d,4f,4d,4f$, summing to
$8(d+f)\leq4(a+d+f)$. Every feasible saturated face is one of
$z_2=\pm d$ or $z_3=\pm f$ and is covered by a displayed edge.
This proves for every choice of nonnegative widths a closed physical
tour covering the saturated section faces with length at most
\begin{equation}\label{eq:tripleperimeter}
 4(a+d+f)=2c\leq4c=2-4b.
\end{equation}
The positive widths in the present allocation ensure that our selected
face points are on its true boundary. The formulas themselves continue
to cover zero and repeated-edge degeneracies.

\subsection{Comparing the same seven-point cycle}

A closed polygonal tour of physical length $P$ admits a parametrization
$\gamma:[0,P]\longrightarrow\R^5$ with $\gamma(0)=\gamma(P)$ and
\[
 \norm{\gamma(t)-\gamma(s)}\leq t-s\qquad(0\leq s\leq t\leq P).
\]
This follows by parametrizing each nonzero straight edge by its physical
length and keeping zero edges constant. Choose a covering parameter for
each of the seven $z_j$ and sort these parameters. The distances between
successive points are bounded by the successive parameter gaps; the
closing distance is at most the sum of the two gaps through the common
endpoint of the closed tour. Hence, in this particular cyclic order,
\[
 \sum_j\norm{z_j-z_{j+1}}\leq P.
\]
Repeated parameters cause no gap in this argument; in the present
application they are also ruled out by
\eqref{eq:quarterseparation}.

Both allocations with five positive slots supply such a tour with
$P\leq2-4b$, by \eqref{eq:doubleperimeter} or
\eqref{eq:tripleperimeter}. Apply the lower bound
\eqref{eq:cyclelower} to exactly this cyclic ordering of the same
seven original labels. It gives
\[
 \frac72-40b\leq2-4b,
 \qquad b\geq\frac1{24},
\]
contrary to our assumption. This completes the proof of
Proposition~\ref{prop:quartetbound}, using
\eqref{eq:quartetwidth} for its width conclusion.

\section{The final five-edge matching}\label{sec:finalmatching}

We have excluded all cases of \eqref{eq:cases} except $(5,5)$. Relabel
the five actual coordinates so that
\[
 L_i=\frac12+t_i,\qquad
 \frac13\geq t_1\geq t_2\geq t_3\geq t_4\geq t_5>0.
\]
Their ten extreme labels $a_i,b_i$ are all distinct, and let $o$ be
the eleventh label. Select the original quartet
\[
 Q=\{o,a_5,b_5,a_4\}.
\]
Its width in coordinate five is $L_5$. In coordinate four it contains
the actual minimum $a_4$, while the other three labels are nonextreme;
\eqref{eq:inner} therefore bounds this coordinate width by $1/2$.
In each of coordinates one, two, and three, all four labels avoid the
two extremes, so their width is at most $1-L_i$. Applying
Proposition~\ref{prop:quartetbound} to this same original quartet gives
\begin{equation}\label{eq:selectedwidth}
 \frac{37}{24}\leq W_Q
 \leq L_5+\frac12+\sum_{i=1}^3(1-L_i)
 =\frac72-W+2L_5+L_4.
\end{equation}
Its seven outsiders are exactly the seven original labels used by the
quartet proof; no new auxiliary equilateral configuration is introduced.

For the actual nonnegative cap slacks $d_i=1/3-t_i$, $1\leq i\leq4$,
inequality \eqref{eq:selectedwidth} is equivalent to
\begin{equation}\label{eq:slack}
 d_1+d_2+d_3+t_5\geq\frac1{24}.
\end{equation}
Indeed it says $t_1+t_2+t_3-t_5\leq23/24$. The fourth slack cancels
in this conversion, and the fifth excess has the positive sign displayed
in \eqref{eq:slack}.

In each wide source row, the other four edges contribute eight distinct
endpoint labels. Apply \eqref{eq:pairedrearrangement} to those four
weights in their inherited descending order and sum the five rows.
The result is
\begin{equation}\label{eq:finallower}
 \Delta\geq F:=\sum_{i=1}^5B_it_i,
\end{equation}
with exactly the coefficients of \eqref{eq:Bdefinitions}. To check the
ordered count explicitly, the coefficient rows, with columns indexed by
$t_1,\ldots,t_5$, are
\begin{equation}\label{eq:coefficientmatrix}
 \begin{pmatrix}
 0&\kappa_1&\kappa_2&\kappa_3&\kappa_4\\
 \kappa_1&0&\kappa_2&\kappa_3&\kappa_4\\
 \kappa_1&\kappa_2&0&\kappa_3&\kappa_4\\
 \kappa_1&\kappa_2&\kappa_3&0&\kappa_4\\
 \kappa_1&\kappa_2&\kappa_3&\kappa_4&0
 \end{pmatrix}.
\end{equation}
For example $t_3$ receives two copies of $\kappa_2$ and two of
$\kappa_3$, and $t_5$ receives four copies of $\kappa_4$.
The two endpoint occurrences per edge are already included in each
$\kappa$ through \eqref{eq:minimalsums}; no further factor is inserted.

All five coordinates are wide, so the exact same energy is
\[
 \Delta=5\alpha-\frac{15}{4}+\frac12\sum_{i=1}^5t_i.
\]
Substitute $t_i=1/3-d_i$ for $i\leq4$ and subtract. This gives the
identity
\begin{equation}\label{eq:slackidentity}
 F-\Delta=G_4+\sum_{i=1}^4\left(\frac12-B_i\right)d_i
                 +\left(B_5-\frac12\right)t_5,
\end{equation}
where
\begin{equation}\label{eq:G4}
 G_4:=\frac{B_1+B_2+B_3+B_4}{3}-5\alpha+\frac{37}{12}
 =\frac{-2851+180\sqrt{37}+60\sqrt{185}+420\sqrt5}{132}.
\end{equation}
The constant $G_4$ need not be positive. The original quartet coupling
\eqref{eq:slack}, rather than its sign alone, supplies the contradiction.

Put $\gamma=1/2-B_3$. By \eqref{eq:Bsigns},
\[
 \gamma>\frac{15}{64},\qquad
 \frac12-B_i\geq\gamma\ (i=1,2,3),\qquad
 \frac12-B_4>0,\qquad B_5-\frac12>\gamma.
\]
Since all the slacks and $t_5$ are nonnegative,
\begin{equation}\label{eq:gapfromslack}
 F-\Delta\geq G_4+\gamma(d_1+d_2+d_3+t_5)
       \geq G_4+\frac\gamma{24}.
\end{equation}
The exact radical certificate in Section~\ref{sec:arithmetic} is
$G_4>-4349/660000$. Therefore
\begin{equation}\label{eq:finalmargin}
 F-\Delta>
 -\frac{4349}{660000}+\frac{15}{64\cdot24}
 =-\frac{4349}{660000}+\frac5{512}
 =\frac{33541}{10560000}>0.
\end{equation}
This contradicts \eqref{eq:finallower}. The five-edge matching is
impossible, and the complete forest classification proves that eleven
original unit-equilateral points in $\R^5$ cannot exist.

\section{Exact arithmetic certificates}\label{sec:arithmetic}

This section supplies the numerical signs used in the proof. Every
comparison concerns nonnegative real numbers and follows by squaring
rational bounds. No decimal approximation is required.

\subsection{The three profile bounds}

For the modified profile, the following integers certify strict lower
square-root bounds: in each row $m^2<Nd^2$, hence $\sqrt N>m/d$.
\begin{center}
\begin{tabular}{rrrr}
\toprule
$N$ & $m$ & $d$ & $Nd^2-m^2$\\
\midrule
$14$ & $37416$ & $10000$ & $42944$\\
$91$ & $95393$ & $10000$ & $175551$\\
$221$ & $148660$ & $10000$ & $204400$\\
$323$ & $179722$ & $10000$ & $2716$\\
\bottomrule
\end{tabular}
\end{center}
Their weighted numerator is
\[
 6(37416)+4(95393)+4(148660)+4(179722)=1919596.
\]
It follows that
\[
 \beta_M=\frac{215-6\sqrt{14}-4\sqrt{91}-4\sqrt{221}-4\sqrt{323}}{22}
  <\frac{57601}{55000}<\frac{21}{20},
\]
as asserted in \eqref{eq:modifiedcap}.

For the nonwide profile,
\[
 139^2<31\cdot25^2,\qquad 1623^2<1054\cdot50^2.
\]
Thus $20\sqrt{31}+4\sqrt{1054}>6026/25$, and
\[
 a_N+\frac13
   =\frac{271-20\sqrt{31}-4\sqrt{1054}}{33}
   <\frac{749}{825}<\frac{10}{11}.
\]

For the unshifted root profile, the following positive differences give
the analogous lower bounds with denominator $1000$:
\begin{center}
\begin{tabular}{rrr}
\toprule
$N$ & $m$ & $N\cdot1000^2-m^2$\\
\midrule
$5$ & $2236$ & $304$\\
$3$ & $1732$ & $176$\\
$42$ & $6480$ & $9600$\\
$210$ & $14491$ & $10919$\\
\bottomrule
\end{tabular}
\end{center}
The weighted sum is
$12(2236)+24(1732)+8(6480)+4(14491)=178204$.
Consequently
\[
 \alpha_0<\frac{2949}{2750}<\frac{43}{40},
\]
which proves \eqref{eq:alpha0} and the strict isolated-mass bound.

\subsection{Positive increments and their upper and lower bounds}

The definitions \eqref{eq:sdefinitions} also have the useful forms
\begin{equation}\label{eq:incrementsquares}
 s_3=\frac{(\sqrt{37}-5)^2}{22},\qquad
 s_4=\frac{(3\sqrt5-\sqrt{37})^2}{22},\qquad
 s_5=\frac{(7-3\sqrt5)^2}{22}.
\end{equation}
The unsquared increments are positive. For the first, $\sqrt{37}>6$.
For the second, the rational square comparisons
\[
 \sqrt5>\frac{559}{250},\qquad
 \sqrt{37}<\frac{6083}{1000}
\]
give $3\sqrt5-\sqrt{37}>5/8$. For the third,
$\sqrt5<7/3$. We may therefore square the following upper bounds
without reversing a sign:
\[
 \sqrt{37}-5<\frac{11}{10},\qquad
 3\sqrt5-\sqrt{37}<\frac34,\qquad
 7-3\sqrt5<\frac25.
\]
They follow respectively from $\sqrt{37}<61/10$, from
$\sqrt{37}>6$ and $\sqrt5<9/4$, and from $\sqrt5>11/5$.
Equation~\eqref{eq:incrementsquares} yields
\begin{equation}\label{eq:supper}
 s_3<\frac{121}{2200}<\frac1{16},\qquad
 s_4<\frac9{352}<\frac1{32},\qquad
 s_5<\frac2{275}<\frac1{128}.
\end{equation}
Similarly $\sqrt{37}>121/20$ gives
$\sqrt{37}-5>21/20$, and the second increment is greater than $5/8$.
Thus
\begin{equation}\label{eq:slower}
 s_3>\frac{441}{8800}>\frac1{20},\qquad
 s_4>\frac{25}{1408}>\frac1{60},\qquad s_5>0.
\end{equation}
All auxiliary root comparisons above are checked by the following
integer inequalities, in addition to $36<37$:
\[
 \begin{gathered}
 559^2=312481<5\cdot250^2=312500,\qquad
 6083^2=37002889>37\cdot1000^2=37000000,\\
 7^2>5\cdot3^2,\quad 61^2>37\cdot10^2,\quad
 9^2>5\cdot4^2,\quad11^2<5\cdot5^2,\quad
 121^2<37\cdot20^2.
 \end{gathered}
\]
The upper and lower bounds imply
$s_1>s_2>s_3>s_4>s_5>0$. They also prove
\eqref{eq:pathsigns}: the first coefficient is less than
$4(1/16+1/32)-1/2=-1/8$, and the second is greater than
$9/22+1/20+3/60-1/2=1/110$.

\subsection{The shifted trace bound and matching constants}

The three root bounds needed for the remaining exact constants are
\begin{equation}\label{eq:finerooots}
 \sqrt{37}>\frac{608276}{100000},\qquad
 \sqrt{185}>\frac{1360147}{100000},\qquad
 \sqrt5>\frac{223606}{100000}.
\end{equation}
Their square comparisons are
\[
 \begin{aligned}
 608276^2&=369999692176<370000000000,\\
 1360147^2&=1849999861609<1850000000000,\\
 223606^2&=49999643236<50000000000.
 \end{aligned}
\]
For $\alpha$, their weighted numerator is
\[
 10(608276)+6(1360147)+42(223606)=23635094,
 \qquad3(23635094)=70905282>70900000.
\]
Thus the subtracted radical sum in \eqref{eq:alpha} exceeds $709/3$,
giving $\alpha<2/3$.

The successive differences of the coefficients in
\eqref{eq:Bdefinitions} are
\[
 \begin{split}
 B_2-B_1&=3s_4,\\
 B_3-B_2&=2s_3+s_4-2s_5,\\
 B_4-B_3&=s_2+2s_3-s_4-2s_5,\\
 B_5-B_4&=3(s_2-s_4).
 \end{split}
\]
They are all positive by the decreasing positive $s_i$ sequence.
Also \eqref{eq:supper} gives
\[
 B_3=2(s_3+2s_4+s_5)<\frac{17}{64},\qquad
 B_4=s_2+4s_3+3s_4<\frac{4211}{8800}<\frac12.
\]
Using \eqref{eq:slower},
\[
 B_5+B_3=4s_2+6s_3+4s_4+2s_5
   >\frac8{11}+\frac3{10}=\frac{113}{110}>1.
\]
This proves all signs in \eqref{eq:Bsigns} and
$\gamma=1/2-B_3>15/64$.

Expanding $G_1$ from \eqref{eq:G1} using the $s_i$ definitions gives
its displayed radical expression. The bounds \eqref{eq:finerooots}
then imply
\[
 \begin{split}
 260(608276)+96(1360147)+420(223606)&=382640392,\\
 G_1&>\frac{382640392-382500000}{13200000}
      =\frac{17549}{1650000}>0.
 \end{split}
\]
The same direct expansion gives \eqref{eq:G4}. Its lower certificate is
\[
 \begin{split}
 180(608276)+60(1360147)+420(223606)&=285013020,\\
 G_4&>\frac{285013020-285100000}{13200000}
      =-\frac{4349}{660000}.
 \end{split}
\]
Finally the common-denominator calculation in
\eqref{eq:finalmargin} has numerator
$-4349\cdot16+5\cdot20625=-69584+103125=33541$ and denominator
$10560000$. This completes all arithmetic certificates.

\section{Positive distances, arbitrary sets, and the witness}\label{sec:lifting}

We finish the proof of Theorem~\ref{thm:main}. If eleven distinct
points had common distance $\delta>0$, replacing every coordinate by
$x_k/\delta$ would give eleven unit-equilateral points, since
\[
 \sum_{k=1}^5\abs{\frac{x_k}{\delta}-\frac{y_k}{\delta}}
   =\frac1\delta\sum_{k=1}^5\abs{x_k-y_k}.
\]
This is impossible by the preceding sections. Any finite equilateral
set with more than ten elements contains eleven distinct elements.
An infinite set also contains eleven distinct elements: choose
successively a point outside the previously chosen finite subset.
Thus an arbitrary equilateral set is finite and has at most ten
elements. This argument places no finiteness assumption on the original
set.

For the lower bound take
\[
 S_{10}=\{\tfrac12 e_i,-\tfrac12 e_i:1\leq i\leq5\}.
\]
These are ten distinct points. Opposite points on the same coordinate
axis have distance one. Points on different axes have two nonzero
coordinate differences of magnitude $1/2$, again giving distance one.
This proves the asserted witness and $e(\ell_1^5)=10$.

\section{Formalization and declaration correspondence}\label{sec:formalization}

The accompanying formal development uses Lean and mathlib
\cite{Lean2021,Mathlib2020}. Points are functions from a finite set of
coordinate indices to $\R$, and their distance is defined literally as
the finite sum of absolute coordinate differences. Its eleven-point
input is an arbitrary function $\{1,\ldots,11\}\to\R^5$.
The unit-equilateral predicate quantifies over every pair of distinct
labels, so it includes all $55$ unordered pairs.

The matrix presentation in Section~\ref{sec:cuts} corresponds to scalar
cut energies in the formal sources: on a zero-sum vector, a scalar
column energy is the sum of its nonnegative gaps multiplied by squared
prefix sums. The formal total-energy and total-trace identities are
the identities \eqref{eq:operator}--\eqref{eq:profile} evaluated in
this form. Thus the matrix notation in the paper abbreviates exactly
the ordered scalar computation used by the formal proof.

The following table identifies the main correspondence. All names are
in the namespace \lean{Kusner5}; entries with a further namespace retain
it in the table.
\begingroup\small
\setlength{\tabcolsep}{4pt}
\begin{longtable}{>{\raggedright\arraybackslash}p{0.43\textwidth}>{\raggedright\arraybackslash}p{0.52\textwidth}}
\toprule
Mathematical statement & Formal declaration\\
\midrule
\endfirsthead
\toprule
Mathematical statement & Formal declaration\\
\midrule
\endhead
\bottomrule
\endfoot
Exact all-pair energy and trace identities &
\lean{unit_equilateral_actual_totalEnergy};
\lean{unit_equilateral_five_totalTrace}\\[3pt]
The physical range bound, Proposition~\ref{prop:range} &
\lean{actual_coordinate_range_le_five_sixths}\\[3pt]
At least four wide coordinates and exact shifted cross energy &
\lean{Profiles.at_least_four_wide};
\lean{Profiles.shifted_crossEnergy}\\[3pt]
At least two isolated edges and all five possible count cases &
\lean{ForestMatching.original_at_least_two_isolated};
\lean{ForestMatching.original_exhaustive_counts}\\[3pt]
The two cases with exactly two isolated edges &
\lean{ForestMatching.no_four_wide_two_isolated};
\lean{ForestMatching.no_five_wide_two_isolated}\\[3pt]
Four-edge matching and three-isolated-edge path &
\lean{FourMatchingAssembly.no_four_matching};
\lean{PathAssembly.no_five_wide_three_isolated}\\[3pt]
Quartet cut totals, original overlap, and all four cross equations &
\lean{quartet_equal_cut_totals};
\lean{quartet_original_cooccurrence};
\lean{quartet_original_cross_equation}\\[3pt]
Actual projection and positive active-face projection &
\lean{exists_original_quartet_projection};
\lean{exists_original_quartet_face_projection}\\[3pt]
Physical slot classification and the same seven-label lower cycle &
\lean{original_quartet_positive_slot_classification};
\lean{original_projected_seven_cycle_lower}\\[3pt]
The two physical boundary tours and cyclic upper bound &
\lean{exists_original_113_section_tour};
\lean{exists_original_122_section_tour};
\lean{QuartetGeometry.seven_cycle_of_closed_path}\\[3pt]
Arbitrary original quartet width, Proposition~\ref{prop:quartetbound} &
\lean{original_quartet_singleton_ge_one_twenty_fourth};
\lean{original_quartet_width_ge_thirty_seven_twenty_fourths}\\[3pt]
Five actual matching rows and the selected quartet slack &
\lean{MatchingAssembly.original_five_matching_energy_lower};
\lean{MatchingAssembly.matching_original_slack}\\[3pt]
Eleven-point unit-distance exclusion, with the literal expanded input &
\lean{no_eleven_unit}; \lean{no_eleven_unit_explicit}\\[3pt]
Positive-distance normalization and finite-array bound &
\lean{no_eleven_positive_distance};
\lean{equilateral_array_card_le_ten}\\[3pt]
Arbitrary-set finiteness and cardinality bound &
\lean{equilateral_set_finite_card_le_ten}\\[3pt]
The ten-point witness and complete exact result &
\lean{exists_equilateral_ten}; \lean{kusner_five_exact}\\
\end{longtable}
\endgroup

In particular, the fully expanded eleven-point theorem has the
mathematical type
\[
 \forall x:\{1,\ldots,11\}\to\R^5,\quad
 \neg\left(\forall i\ne j,\ \sum_{k=1}^5\abs{x_{i,k}-x_{j,k}}=1\right).
\]
The final exact theorem combines the arbitrary-set, arbitrary-positive-
distance upper bound with the explicit finite ten-point lower witness.
The arbitrary-quartet lower bound is proved from the original
configuration before it is supplied to the final matching argument;
it is not an additional hypothesis of the final theorem.

The development reuses mathlib's real-number, finite-sum, finite-set,
sorting, and square-root results. The project declarations and their
composition are kernel-checked, and the transitive axiom inspection of
the central and final declarations yields only the standard foundations
\lean{propext}, \lean{Classical.choice}, and \lean{Quot.sound}.
All numerical certificates above are represented by exact arithmetic
proofs. No unproved mathematical axiom or admitted proof is used.
This verification statement concerns the supplied formal development
with its declared library dependencies.

\section*{Acknowledgement}
Large language model (LLM)-based tools assisted with the research and preparation of this manuscript.

\bibliographystyle{plain}
\bibliography{references}

\end{document}
